\documentclass[11pt,a4paper]{article}
\usepackage[T1]{fontenc}
\usepackage{lmodern,amsmath,amssymb,amsthm,mathtools,booktabs,microtype}
\usepackage[margin=25mm]{geometry}
\usepackage[colorlinks=true,linkcolor=blue,citecolor=blue,urlcolor=blue]{hyperref}
\setlength{\emergencystretch}{2em}
\newtheorem{theorem}{Theorem}[section]
\newtheorem{lemma}[theorem]{Lemma}
\newtheorem{proposition}[theorem]{Proposition}
\newtheorem{corollary}[theorem]{Corollary}
\newcommand{\Q}{\mathbb Q}\newcommand{\Z}{\mathbb Z}
\DeclareMathOperator{\Tr}{Tr}\DeclareMathOperator{\Nm}{N}
\title{Small polynomial classes for the actual detector trace}
\author{}\date{17 November 2026}
\begin{document}\maketitle
\begin{abstract}
Three polynomial square classes preserve the trace pole test without
losing a contact prime or introducing a rational coordinate. Their
square radicands have coefficient height four in the trace and weight;
the smallest is a cubic with four monomials. We prove compatibility
with every selected contact and fixed preparation place, retain all
whole-character guards, and give the associated curve models. The smallest
cubic has a rational point of infinite order for every rational normalized
weight $0<w\le1/4$. The
resulting construction criteria still require rational realization at
the actual detector traces.
\end{abstract}

\section{The three square tests}

Let $S$ be the retained fixed exceptional set, with $2\in S$, and
let $w,V\in\Q^\times$. Assume that every odd valuation of $w$
outside $S$ is positive. For the native elliptic sources,
\[
 V=\Tr_{K/\Q}(W^{-1}),\qquad K=\Q(\sqrt2),\qquad
 b=4w/V^2,
\]
with the normalized general or conditional physical weight of
\cite{Reweighted,Trace}. Define
\begin{equation}\label{eq:classes}
 \mathcal C=V^2+w(2-V),\qquad
 \delta_1=V+2,\quad
 \delta_2=V^2-V+2,\quad
 \delta_3=1+w(2-V).
\end{equation}
The tests use the actual retained $w,V$; these are not new free
parameters in the source.

\begin{theorem}\label{thm:pole}
For each $i\in\{1,2,3\}$, the condition
\[
                       \delta_i\mathcal C\in\Q^{\times2}
\]
excludes every negative odd valuation of $b$ outside $S$.
At every prepared real character, where $w>0$ and $0<V\le2$,
all three radicands are strictly positive. No contact-prime or
height-support allowance is needed for these classes.
\end{theorem}
\begin{proof}
If $v_p(b)<0$ is odd, put $\nu=v_p(w)$ and $m=v_p(V)$.
Then $\nu>0$ is odd and $2m>\nu$, so $m\ge1$. Since $p\ne2$,
the number $2-V$ is a unit. Thus $v_p(\mathcal C)=\nu$.
The three multipliers reduce, respectively, to $2,2,1$ modulo $p$.
They are units, so each product has odd valuation and cannot be
a rational square. Squareness of these unit residues is not
required for this contradiction.

On the prepared real domain, $\mathcal C\ge V^2>0$,
$\delta_1>2$, and
\[
 \delta_2=(V-\tfrac12)^2+\tfrac74>0,\qquad \delta_3\ge1.
\]
This proves the positivity assertions. The pole implication holds
at every original nonexceptional contact, so the inherited local
extraction and height arguments retain their full prime support.
\end{proof}

\section{The selected primes and local preparation}

\begin{lemma}\label{lem:splitting}
The retained simultaneous prime construction may use precisely its
existing avoidance modulus while ensuring that every selected contact
prime is congruent to one modulo $24$. In particular both two and
three are squares in its residue field. The conjugate reduction-order
conditions and arbitrary fixed denominator divisibility are preserved.
\end{lemma}
\begin{proof}
The arbitrary-avoidance lemma of \cite{Arithmetic} includes three
in the odd avoidance modulus. In the primary source
\cite[Theorem 4.1]{Sun}, the selected Gaussian prime generators
satisfy $\pi\equiv9\pmod{16}$ and $\pi\equiv i\pmod M$;
the allowed odd squarefree modulus $M$, prime to seven, may contain
three. Taking norms gives $p\equiv1\pmod8$ and
$p\equiv1\pmod3$. Therefore $p\equiv1\pmod{24}$.
Quadratic reciprocity gives $(2/p)=(3/p)=1$.
These are conditions of the same theorem, which simultaneously
supplies its cross-order and self-order gcds, finite avoidance, and
prescribed square-denominator dilation. No separate prime existence
assertion is being added.
\end{proof}

\begin{proposition}\label{prop:local}
All three tests in \eqref{eq:classes} are nonzero local squares
at every selected comparison contact of the retained preparation.
They are also nonzero local squares at every fixed preparation place
after taking the original trace preparation to sufficient precision.
\end{proposition}
\begin{proof}
At a selected contact $v_p(w)>0$ and the actual trace has residue
one or two. The resulting residues are
\[
\begin{array}{c|cccc}
V&\mathcal C&\delta_1&\delta_2&\delta_3\\\hline
1&1&3&2&1\\
2&4&4&4&1
\end{array}
\]
and are nonzero squares by Lemma~\ref{lem:splitting}.
Hensel's lemma therefore proves the selected-contact assertion.
The trace residues and split-place preparation are those of
\cite{General,Conditional,Native}; the square conditions here do
not change the auxiliary index congruences.

At a fixed place, the existing source preparation makes $V$
arbitrarily close to two while $w$ stays in a prescribed bounded
neighborhood. In the conditional source $w$ tends to zero.
The three classes approach $4,4,1$, respectively, and
$\mathcal C$ approaches four. These are nonzero rational squares,
whose square classes have open neighborhoods in every
$\Q_p^\times$.

For an explicit check at two, put $e_0=V-2$. Then
\[
 \delta_1=4+e_0,\quad
 \delta_2=4+3e_0+e_0^2,\quad
 \delta_3=1-we_0,\quad
 \mathcal C=V^2-we_0.
\]
It suffices to choose
$v_2(e_0)\ge5$ and $v_2(we_0)\ge5$. The first two classes
divided by four, the third class itself, and
$\mathcal C/V^2$ then lie in $1+8\Z_2$.
The general source first fixes its scalar neighborhoods and then
its slope, so the required precision is available in its existing
order of preparation. At odd fixed places the same proof uses
$1+p\Z_p$. All earlier denominator and point avoidance conditions
are retained.
\end{proof}

These assertions concern the selected contacts and the fixed finite
preparation set. They do not assert local squareness at every
remaining prime, nor rational squareness of the actual radicands.

\section{Exact polynomial blocks}

Put $R_i=\delta_i\mathcal C$. Direct expansion gives
\begin{align}
 R_1&=V^3+(2-w)V^2+4w,\label{eq:smallcubic}\\
 R_2&=V^4-V^3-wV^3+2V^2+3wV^2-4wV+4w,\label{eq:quartic}\\
 R_3&=-wV^3+(w+1)^2V^2-w(4w+1)V+2w(2w+1).
 \label{eq:weightcubic}
\end{align}
The exact statistics as integer polynomials in the two formal
variables $(V,w)$ are
\begin{center}
\begin{tabular}{lrrrr}\toprule
Radicand&Total degree&Degree in $V$&Monomials&Coefficient height\\\midrule
$R_1$&3&3&4&4\\
$R_2$&4&4&7&4\\
$R_3$&4&3&8&4\\\bottomrule
\end{tabular}
\end{center}
The equations $Y^2-R_i=0$ have respectively five, eight and nine
monomials and the same total degrees and coefficient height.
These are block coefficients before substitution of the actual
weight and trace and before determinant compilation. They are not
degree or coefficient-height bounds for a complete defining polynomial.

\section{Curve models and rational reference points}

For the smallest cubic, the substitution $x=V+2$ gives
\begin{equation}\label{eq:modelone}
 Y^2=x\bigl(x^2-(w+4)x+4(w+1)\bigr).
\end{equation}
The reference point is $(V,Y)=(2,4)$, or $(x,Y)=(4,4)$.
The point $(x,Y)=(0,0)$ has order two when the curve is smooth.
The Weierstrass discriminant is
\[
                       256w(w-8)(w+1)^2.
\]
For $R_2$, the reference point is again $(2,4)$, and the
discriminant of the quartic in $V$ is
\[
                       -112w(w-8)(w-1)^4.
\]
Thus both models are smooth at every prepared weight
$0<w\le1/4$.

For $R_3$, the rational change
\[
 x=1+w(2-V),\qquad z=wY
\]
gives
\begin{equation}\label{eq:modelthree}
 z^2=x\bigl(x^2+(w^2-4w-2)x+(w+1)(3w+1)\bigr),
\end{equation}
whose discriminant is
\[
                   16w^3(w-8)(w+1)^2(3w+1)^2.
\]
This model is also smooth for $0<w\le1/4$ and has the rational
point of order two $(x,z)=(0,0)$.

The two rational sections
\[
 Q=(V,Y)=(1,1+w),\qquad R=(2,2)
\]
satisfy $R=-[2]Q$. Indeed, if $f(V)=R_3(V,w)$, then
\[
 f(V)-\bigl((1-w)V+2w\bigr)^2=-w(V-2)(V-1)^2,
\]
so the line through $Q,R$ is tangent at $Q$. Translation by the
two-torsion point sends the trace coordinate to
\[
                V\longmapsto
                \frac{w-(1+2w)V}{1+w(2-V)}.
\]
The displayed coordinates are changes of variables for studying the
square equations, not additional free coordinates in their finite-algebra
encoding. They do not prescribe the actual detector trace. The next
section proves positive rank for the smallest cubic throughout the
normalized range.

\section{Positive rank throughout the normalized weight range}

\begin{proposition}\label{prop:smallcubic-rank}
For every rational $w$ with $0<w\le1/4$, the point $P=(4,4)$
has infinite order on
\[
 E_w:\quad y^2=x\bigl(x^2-(w+4)x+4(w+1)\bigr).
\]
In particular, $E_w(\Q)$ has positive rank.
\end{proposition}
\begin{proof}
The discriminant is $256w(w-8)(w+1)^2$, which is nonzero in the
stated range. The quadratic factor of the cubic has discriminant
$w(w-8)<0$. Thus $(0,0)$ is the unique nonzero rational point of
order two. Mazur's torsion theorem
\cite[Introduction, Theorem~(8), p.~35]{Mazur1977} shows that the
rational torsion subgroup is cyclic of order $2,4,6,8,10$, or $12$.
If $P$ were torsion, its order would therefore divide one of
$8,10,12$.

Write $\Psi_n(w)=\psi_n(P)$ for the standard division polynomials
evaluated at $P$, with $\Psi_1=1$ and $\Psi_2=8$. Their initial
values and recurrences are
\begin{align*}
 \Psi_3&=-16(w-7)(w+1),\\
 \Psi_4&=1024(w-3)(w^2-8w-1),\\
 \Psi_{2m+1}&=\Psi_{m+2}\Psi_m^3-
                     \Psi_{m-1}\Psi_{m+1}^3,\\
 8\Psi_{2m}&=\Psi_m\bigl(\Psi_{m+2}\Psi_{m-1}^2-
                     \Psi_{m-2}\Psi_{m+1}^2\bigr).
\end{align*}
The recurrences give the following exact primitive data. Here
$c_n$ is the positive content of $\Psi_n$ and $Q_n=\Psi_n/c_n$.
\begin{center}
\begin{tabular}{rrrr}\toprule
$n$&$c_n$&$\deg Q_n$&Leading coefficient of $Q_n$\\\midrule
8&$2^{35}$&15&1\\
10&$2^{51}$&24&5\\
12&$2^{74}$&35&3\\\bottomrule
\end{tabular}
\end{center}
All coefficient arrays and their contents are obtained by integer
arithmetic from the displayed recurrences.

By the rational-root theorem, a rational zero of $Q_8$ is an
integer, whereas a rational zero of $Q_{12}$ has reduced denominator
dividing three. Neither polynomial has a rational zero in
$(0,1/4]$. A rational zero of $Q_{10}$ in this interval could only
be $1/5$. Reduction modulo three excludes that value: at $w=2$,
the initial values and recurrences give
\[
 (\Psi_2,\Psi_3,\Psi_4,\Psi_5,\Psi_7,\Psi_{10})
                    \equiv(2,0,1,2,1,1)\pmod3.
\]
Since $1/5\equiv2\pmod3$ and $2^{51}\equiv2\pmod3$, it follows
that $Q_{10}(1/5)\equiv2\pmod3$. Consequently none of
$\psi_8(P),\psi_{10}(P),\psi_{12}(P)$ vanishes. This contradicts
every possible torsion order and proves the proposition.
\end{proof}

The proposition supplies positive rank for each fixed normalized
weight. Realization of the original detector trace
$V=\Tr_{K/\Q}(W^{-1})$ still requires a rational point whose
abscissa is precisely $V+2$; positive rank alone does not prescribe
that abscissa. Thus this result does not establish the simultaneous
trace realization required for the proposed $21$- and $9$-variable
constructions.

\section{Whole-character guards and the construction criteria}

\begin{theorem}\label{thm:encoding}
The three multipliers may be used independently in the two trace
blocks of the retained general or conditional source. They add no
unrestricted rational coordinate, denominator divisor or lost
contact prime. Under the source's retained arithmetic hypotheses,
the resulting equations are sound. If every required integer input
admits the corresponding simultaneous rational squares at its actual
guarded source tuple, the coordinate counts are twenty-one and nine,
respectively. That realization property is an additional hypothesis
and is not proved here.
\end{theorem}
\begin{proof}
At a prepared point tuple every geometric character of the finite
point algebra is real. The two real slots of every $W$ are at least
one, so $w>0$ and $0<V\le2$ there. Theorem~\ref{thm:pole}
therefore makes each selected $R_i$ positive at every character.
The three classes are polynomials in the already guarded $V,w$;
their denominators introduce no new divisor.

Write each chosen radicand in the localized point algebra as $N/D$.
Adjoin the monic equation $Y^2=ND$. Include the numerator and
denominator norms in the scalar guard product, as in \cite{Unit}.
Since $ND=D^2R_i>0$ at every prepared character, these norm factors
are nonzero and discard no prepared tuple. They are scalar guard
factors, not new rational nonzero-witness coordinates. The resulting
quadratic extensions are finite, etale and totally real on that locus.

With $p=9$ point blocks in the general source, or $p=3$ in the
conditional source, the algebra has rank $4^p2^2$. The guarded
separator theorem of \cite{Unit} uses one rational coordinate and
ensures that a rational determinant zero gives a rational character
of the entire algebra. In particular every chosen square root is
rational, not merely real algebraic. The external-coordinate census is
\[
 1\text{ input}+2p\text{ point coordinates}
          +1\text{ common slope}+1\text{ encoder}=2p+3.
\]
There is no extra square, height-recovery or guard coordinate.
Theorem~\ref{thm:pole} supplies the same local pole implication as
the source, so all inherited integer-extraction and height arguments
apply. The stated actual-source realization supplies completeness.
The original conditional input factor includes zero.
\end{proof}

The fixed finite menu $\{1,\delta_1,\delta_2,\delta_3\}$ is
available in each block. Compile each pair of menu choices separately
and multiply the resulting polynomials, identifying all external
coordinates, including the encoder coordinate. A finite product
vanishes over $\Q$ exactly when one branch vanishes. This permits
the class choice to vary with the integer input without an additional
coordinate or an unrestricted class parameter.

The same local and guard proofs apply to a single-cover conditional
source retaining one rational square parameter in place of the common
slope. That source has rank $128$ per branch and the same nine-coordinate
census. Its separate mixed-power height inequality and actual-source
realization requirements remain in force. In every case the established
complete defining counts remain twenty-five and thirteen until the
corresponding rational realization is proved.

\begin{thebibliography}{9}\raggedright
\bibitem{Mazur1977}
B.~Mazur, \emph{Modular curves and the Eisenstein ideal},
Publications Math\'ematiques de l'IH\'ES \textbf{47} (1977),
33--186, Introduction, Theorem~(8), p.~35; III, Theorem~(5.1).
\url{https://www.numdam.org/item/PMIHES_1977__47__33_0.pdf}.

\bibitem{Sun}Z.-W.~Sun,
\emph{48 unknowns for Hilbert's tenth problem over the rationals},
arXiv:2610.20691v3 (2026), Theorem 4.1,
\url{https://arxiv.org/pdf/2610.20691}.
\bibitem{Arithmetic}\emph{Arithmetic dependencies of the construction
with 29 variables}, repository manuscript, arbitrary-avoidance lemma,
\nolinkurl{papers/arithmetic_dependencies/arithmetic_dependencies_20261108_v8.tex}.
\bibitem{General}\emph{A defining polynomial in twenty-five variables:
combined scalar and comparison conditions}, repository manuscript,
\nolinkurl{papers/few_variables25/few_variables25_merged_20261110_v19.tex}.
\bibitem{Conditional}\emph{Thirteen Variables from Direct Elliptic
Index Comparison}, repository manuscript,
\nolinkurl{papers/conditional_few_variables/direct_input13_20261110_v22.tex}.
\bibitem{Reweighted}\emph{Quadratic reweightings of native elliptic
detectors}, repository manuscript,
\nolinkurl{papers/pole_parameters/reweighted_elliptic_detectors_20261116_v53.tex}.
\bibitem{Native}\emph{A guarded rational-square criterion for normalized
elliptic detectors}, repository manuscript, fresh split contacts,
\nolinkurl{papers/pole_parameters/native_square_criterion_20261116_v53.tex}.
\bibitem{Trace}\emph{A rational conic criterion using the actual detector
trace}, repository manuscript,
\nolinkurl{papers/pole_parameters/trace_conic_criterion_20261117_v57.tex}.
\bibitem{Unit}\emph{A single norm determinant for the 25- and 13-variable
constructions}, repository manuscript, guarded rational-character theorem,
\nolinkurl{papers/finite_algebra/unit_norm_20261114_v28.tex}.
\end{thebibliography}
\end{document}
